\section{Discussion and Conclusion}
\label{sec:discussion}

The principal result is a boundary on fixed recovery: it improves timely
confirmation in the disrupted-path setting but worsens it under the
higher offered load. Local deadline feasibility is insufficient to decide
whether a recovery flood is beneficial to a shared channel. The clean
control reveals this limitation even though recovery-associated deliveries
exist and full CPR has higher destination delivery. Contention-aware
admission is a motivated future direction, not a policy evaluated here.

The external behavioral models omit documented retry and learning
behavior; custom frame lengths omit product framing and security overhead.
Parameter matching does not resolve these limits, so the results do not
establish product-firmware superiority. Twenty seeds, three scenarios,
and three deadlines provide finite descriptive evidence, not generality
across topology, traffic, and radio settings.

The radio model includes half-duplex operation, overlapping transmissions,
capture/interference decisions, and probabilistic reception, but has no
hardware calibration in this study. LoRa capture, imperfect spreading-factor
orthogonality, and aggregate interference can change real reception
behavior~\cite{rahmadhani2018concurrent,croce2018orthogonality,heusse2023capacity}.
The grid and fault fields are controlled stress settings. No duty-cycle or
listen-before-talk (LBT) enforcement is modeled. Fixed PHY, local timing,
and simplified queues leave device CAD timing, clock error, and application
diversity outside the claimed scope.
Transmit airtime does not measure whole-device energy or battery life.

Finite source stages and state limits do not guarantee delivery or bound
all simulator/firmware memory. Zero evictions do not measure peak ledger
occupancy. Canceled copies can change later contention, so their counts
alone do not establish network cost savings.

MeshEcho-CPR makes source-stage admission, route-evidence concurrency,
and the relay commitment boundary explicit. In the completed study,
same-path repeat improves confirmation in both tested component settings,
and ACK cancellation modestly reduces airtime. Recovery supplies a
conditional benefit with a substantial failure mode: its S2 gain of
8.32~pp becomes a 15.99-pp loss in S3. A useful recovery design must
therefore account for channel contention as well as remaining time.
